% ID: SEL-D0024-OB-E03
% ESTADO: cotejada. Segunda pasada visual completa y borrador compilado revisados.
% Fuente: Selectividad/matematicas_aplicadas/reserva_b_EXAMEN 2-ANDALUCÍA.pdf
% SHA-256: 98b7723232acd3143de28397f7641923ed48292db7c35883ebd2b56a72b30b79
% Página: 2; opción B; ejercicio original 3.

A una asamblea en la Universidad asisten 420 alumnos de los cuales 180 son de Empresariales, 72 de Relaciones Laborales y el resto de Derecho. Un tercio de los alumnos de Empresariales, dos tercios de los de Derecho y 16 alumnos de Relaciones Laborales votan NO a la huelga. El resto ha votado SÍ.
\begin{enumerate}[label=\alph*)]
\item \textbf{(0.9 puntos)} Calcule la probabilidad de que elegido un alumno al azar, sea de Empresariales y haya votado SÍ a la huelga.
\item \textbf{(0.8 puntos)} \textquestiondown{}Cuál es la probabilidad de que elegido un alumno al azar haya votado SÍ a la huelga?
\item \textbf{(0.8 puntos)} Si elegido un alumno al azar, resulta que ha votado NO a la huelga, \textquestiondown{}cuál es la probabilidad de que sea de Relaciones Laborales?
\end{enumerate}
